\chapter{Closed-Cell Engaged-Null Catalogue}\label{app:engaged-null}

This appendix collects the seven engaged-null entries and one
saturated cluster that comprise the closed-cell evidence referenced
by the synthesis section of Chapter~\ref{ch:experiments}. Each entry
is a mechanism that activated under stochastic gradient descent on
the LibriSpeech causal RWKV-6 spine without converting into a
measurable test CER reduction. Engagement is established via
parameter mobility diagnostics and per-block realised statistics;
the absence of CER reduction is the marker for axis-mismatch with
the LibriSpeech task structural prior. The catalogue and its
adjacent saturated cluster form the converse evidence for the
cross-experiment invariant articulated in
Chapter~\ref{ch:experiments}: dense per-token freedom does not
convert into measured gain without alignment between the
mechanism's function-class extension and the axis the task
exercises.

The empirical numbers in this appendix are taken from the
discovery-phase per-stage records in
\texttt{experiments/formal\_v1/} (the
\texttt{stages\_2\_9\_summary.md},
\texttt{STAGE10\_SUMMARY.md}, and
\texttt{EXPRESSIVITY\_AXES.md} documents); the runs were trained at
$7$~M parameters, $30$ epochs, on LibriSpeech clean-100, seed 42,
with the noise floor on the discovery-phase codebase estimated at
$\sigma \approx 0.0014$. The vanilla causal RWKV-6 reference for
the discovery phase is dev CER $0.1258$ / test CER $0.1263$.

\section{Seven engaged-null mechanisms}\label{app:engaged-null:seven}

\subsection{Delta rule rank-one Householder erasure}\label{app:engaged-null:delta}

The delta rule mechanism replaces the per-token recurrent update with
a rank-one Householder erase
$S_{t} = w_{t} \odot (I - \beta_{t} k_{t} k_{t}^{\top}) S_{t-1} +
k_{t} v_{t}^{\top}$, with the gate $\beta_{t}$ data-dependent through
a sigmoid head and a per-head scalar
$g_{\delta}$~\cite{schlag2021deltanet, yang2024deltanet}. The
mechanism targets axis 2: the rank-one erasure subtracts the state
component aligned with the current key before writing the new value,
reducing readout interference under associative-memory load.

\begin{table}[h]
\centering
\small
\caption[Delta rule engaged-null entry]{Delta rule fact box.}
\label{tab:engaged-null:delta}
\setlength{\tabcolsep}{6pt}
\renewcommand{\arraystretch}{1.05}
\begin{tabularx}{\linewidth}{@{}p{3.3cm} X@{}}
\toprule
Target axis & 2 (associative memory under interference) \\
Cell & \texttt{rwkv6\_delta\_warmstart\_fixed} \\
Dev / Test CER & $0.1258$ / $0.1256$ \\
$\Delta$ test vs vanilla & $-0.0007$ (within $\sigma$) \\
Mobility evidence & Per-head $g_{\delta}$ moves from zero-init to
$0.42$ to $0.65$ across layers; $\beta_{t}$ p95 reaches $1.48$;
realised state-norm erasure up to $24\%$ per token at layer~5 \\
Engaged-null reading & SGD finds and uses the delta-rule
freedom (state is materially modified per token), but LibriSpeech
sequences sit far below the state-saturation regime ($T \le 500$,
$d^{2} = 4096$); the erasure decorrelates the state, but the task
signal does not require axis-2 capacity at this scale, so the
mechanism's modification of the state does not translate into CER
reduction \\
\bottomrule
\end{tabularx}
\end{table}

\subsection{NCGRU-Cayley orthogonal transition}\label{app:engaged-null:cayley}

The NCGRU-Cayley mechanism replaces the diagonal transition with a
per-token Cayley-orthogonal operator
$(I - A_{t})(I + A_{t})^{-1}$ where $A_{t}$ is skew-symmetric and
data-dependent~\cite{mucllari2025crt}. The mechanism targets axes
2 and 3 simultaneously: the orthogonal transition preserves state
norm exactly while introducing a per-token rotation freedom that, at
sufficient depth, supports finite-automaton simulation.

\begin{table}[h]
\centering
\small
\caption[NCGRU-Cayley engaged-null entry]{NCGRU-Cayley fact box.}
\label{tab:engaged-null:cayley}
\setlength{\tabcolsep}{6pt}
\renewcommand{\arraystretch}{1.05}
\begin{tabularx}{\linewidth}{@{}p{3.3cm} X@{}}
\toprule
Target axis & 2 + 3 (state-tracking-adjacent rotation) \\
Cell & \texttt{rwkv6\_orthogonal} \\
Dev / Test CER (epoch 15) & $0.1518$ / not completed (regression
trajectory) \\
$\Delta$ test vs vanilla & $+0.0255$ at epoch~15, regression-track
projection beyond vanilla \\
Mobility evidence & Skew-symmetric $A_{t}$ carries non-zero norm
across all layers; the Cayley operator is exercised but cannot be
absorbed by the existing projections \\
Engaged-null reading & The rotation freedom is dense per-token and
not aligned with the LibriSpeech short-range axis-1 prior; the
norm-preserving constraint additionally interacts adversely with
CTC's character-window decoding at this depth \\
\bottomrule
\end{tabularx}
\end{table}

\subsection{M$^{2}$RNN non-linear state transition}\label{app:engaged-null:m2rnn}

The M$^{2}$RNN mechanism replaces the linear state update with
$S_{t} = \tanh(S_{t-1} W + k_{t} v_{t}^{\top})$ at a single
selected layer~\cite{mishra2026m2rnn}, with a scalar per-head forget
gate. The mechanism targets axis 3: the non-linear state transition
is the only member of its mechanism family in the discovery-phase
catalogue, and crosses into the $\mathsf{NC}^{1}$ class boundary that
diagonal SSMs cannot reach by themselves.

\begin{table}[h]
\centering
\small
\caption[M$^{2}$RNN engaged-null entry]{M$^{2}$RNN fact box.}
\label{tab:engaged-null:m2rnn}
\setlength{\tabcolsep}{6pt}
\renewcommand{\arraystretch}{1.05}
\begin{tabularx}{\linewidth}{@{}p{3.3cm} X@{}}
\toprule
Target axis & 3 (non-Abelian state-tracking) \\
Cell & \texttt{rwkv6\_m2rnn\_sparse} (sparing-use at layer~5) \\
Dev / Test CER & $0.1276$ / $0.1264$ \\
$\Delta$ test vs vanilla & $+0.0001$ (tied within $\sigma$) \\
Mobility evidence & $\tanh$ saturation rate moves above the linear
regime at the deployed layer; the selected layer's state-norm
distribution shifts under training \\
Engaged-null reading & Axis-3 capacity is exercised by tasks
involving non-Abelian state-tracking, which LibriSpeech CTC at $T
\le 500$ does not require; the mechanism's natural test bed is
synthetic permutation-composition or Dyck-$k$ ($k \ge 2$) tasks,
which are deferred to future work as discussed in
Section~\ref{sec:bounds} \\
\bottomrule
\end{tabularx}
\end{table}

\subsection{Non-normal RSE polar parameterisation
(T2)}\label{app:engaged-null:t2}

The non-normal RSE mechanism extends the block-complex transition of
Section~\ref{sec:primitives:dho} with a symmetric anisotropy factor
$P(\rho, \psi) = R(\psi)^{\top} \mathrm{diag}(e^{\rho}, e^{-\rho})
R(\psi)$ on each $2 \times 2$ block, where $\rho$ is per-token and
zero-initialised. The mechanism targets axes 2 and 3 simultaneously:
the non-normal extension introduces dense per-token transition-side
freedom that goes strictly beyond the normal RSE family.

\begin{table}[h]
\centering
\small
\caption[Non-normal RSE T2 engaged-null entry]{Non-normal RSE T2
fact box.}
\label{tab:engaged-null:t2}
\setlength{\tabcolsep}{6pt}
\renewcommand{\arraystretch}{1.05}
\begin{tabularx}{\linewidth}{@{}p{3.3cm} X@{}}
\toprule
Target axis & 2 + 3 (dense per-token polar transition) \\
Cell & \texttt{rwkv6\_nonnormal\_rse\_viscosity} \\
Dev / Test CER & $0.1202$ / $0.1200$ \\
$\Delta$ test vs anchor & $+0.0023$ vs the depth-graded RSE anchor
$0.1177$ (within $\sigma$) \\
Mobility evidence & $\rho$-LoRA Frobenius norm reaches $18$ to $30$
per layer; realised $\rho$ concentrates at the edge layers (mean
$|\rho| = 0.18$ at L0 and $0.11$ at L5); per-block non-normality
score reaches $0.10$ at edge layers; cross-head bimodal
specialisation pattern \\
Engaged-null reading & The dense per-token polar freedom is
substantially used by SGD (mobility is among the largest in the
catalogue), but the LibriSpeech axis-1 task structure does not
reward the additional non-normal directions: the mechanism shifts
the operator family but does not align with the prior \\
\bottomrule
\end{tabularx}
\end{table}

\subsection{Sparse non-normal RSE edge-only
(S9B)}\label{app:engaged-null:s9b}

The sparse non-normal RSE edge-only variant restricts the polar
transition to the first and last layers and runs plain block-complex
RSE in the middle layers. The mechanism targets axes 2 and 3 in the
same sense as T2, with the structural restriction matched to the
T2 diagnostic finding of edge-layer concentration.

\begin{table}[h]
\centering
\small
\caption[Sparse non-normal RSE engaged-null entry]{Sparse non-normal
RSE edge-only fact box.}
\label{tab:engaged-null:s9b}
\setlength{\tabcolsep}{6pt}
\renewcommand{\arraystretch}{1.05}
\begin{tabularx}{\linewidth}{@{}p{3.3cm} X@{}}
\toprule
Target axis & 2 + 3 (sparse edge-restricted polar transition) \\
Cell & \texttt{rwkv6\_sparse\_nonnormal\_rse\_edge\_only\_viscosity} \\
Dev / Test CER & $0.1218$ / $0.1216$ \\
$\Delta$ test vs anchor & $+0.0039$ vs the depth-graded RSE anchor
$0.1177$ \\
Mobility evidence & Per-layer gates settle near $0.5$ to $0.7$ across
the active edge layers; $\rho$ on the gated edge layers behaves
qualitatively like T2's edge-layer realisation; spectral radius max
$\sim 1.0002$ (regression-stability borderline) \\
Engaged-null reading & The mechanism reproduces the T2 engagement
under a structural sparsity that matches the T2 diagnostic; the
sparser deployment did not align the freedom with the axis-1
LibriSpeech prior any better than the dense T2 deployment \\
\bottomrule
\end{tabularx}
\end{table}

\subsection{Readout gauge A1$'$}\label{app:engaged-null:a1prime}

The readout gauge A1$'$ mechanism introduces a data-dependent phase
on the readout
$y_{t} = \mathrm{Re}\!\left( \sum_{b} e^{-i \phi_{t, h, b}}\,
\overline{r_{c, b}}\, S_{t, b} \right)$ with
$\phi_{t, h, b} = W_{\phi} x_{t}$, zero-initialised. The mechanism
targets axis 5: the per-token readout phase introduces a
content-adaptive gauge alignment between the rotated latent state
and the axis-aligned observation vector.

\begin{table}[h]
\centering
\small
\caption[Readout gauge A1$'$ engaged-null entry]{Readout gauge A1$'$
fact box.}
\label{tab:engaged-null:a1prime}
\setlength{\tabcolsep}{6pt}
\renewcommand{\arraystretch}{1.05}
\begin{tabularx}{\linewidth}{@{}p{3.3cm} X@{}}
\toprule
Target axis & 5 (content-adaptive readout phase) \\
Cell & \texttt{rwkv6\_rse\_dphi\_viscosity} \\
Dev / Test CER & $0.1217$ / $0.1207$ \\
$\Delta$ test vs anchor & $+0.0030$ vs the depth-graded RSE anchor
$0.1177$ (regression edge) \\
Mobility evidence & $W_{\phi}$ Frobenius norm reaches $5.5$ to $6.4$
per layer; realised $|\phi|$ p99 $\approx 2$~rad; mobility is among
the largest in the catalogue \\
Engaged-null reading & The mechanism is engaged but acts adversely:
the data-dependent phase pushes more content into the discarded
imaginary axis at the deeper layers, breaking the implicit gauge
alignment that the anchor model already exploits as
cross-block cancellation; axis-5 readout adaptivity is
exercised by tasks with content-driven cross-token routing
demands, which LibriSpeech does not present \\
\bottomrule
\end{tabularx}
\end{table}

\subsection{Content-conditional $\alpha_{d}$ (CB-3)}\label{app:engaged-null:cb3}

The content-conditional $\alpha_{d}$ mechanism replaces the per-layer
mixing weights of MSDC with a per-token gating function
$\alpha_{d, t}^{(\ell)} = \sigma(W_{g}^{(\ell)} x_{t})$, with the
weight $W_{g}^{(\ell)}$ initialised to bias the gate towards the
per-layer mixture used by MSDC. The mechanism targets axis 5: it
adds per-token receptive-field selection to the input-side
multi-scale convolution.

\begin{table}[h]
\centering
\small
\caption[Content-conditional $\alpha_{d}$ engaged-null
entry]{Content-conditional $\alpha_{d}$ fact box.}
\label{tab:engaged-null:cb3}
\setlength{\tabcolsep}{6pt}
\renewcommand{\arraystretch}{1.05}
\begin{tabularx}{\linewidth}{@{}p{3.3cm} X@{}}
\toprule
Target axis & 5 (content-adaptive receptive-field selection) \\
Cell & \texttt{rwkv6\_convshift\_multidil\_symmetric\_gated\_v2} \\
Dev / Test CER & $0.1150$ / $0.1136$ \\
$\Delta$ test vs MSDC alone & $+0.0136$ vs MSDC at $0.1000$
(regression band) \\
Mobility evidence & Per-token gate values move across the unit
interval; per-layer gate distribution shifts under training; the
mechanism is structurally distinct from the per-layer constant
$\alpha_{d}$ baseline \\
Engaged-null reading & Per-token receptive-field selection is dense
per-token freedom on top of a working axis-1 base; LibriSpeech
phoneme-to-syllable acoustic structure does not reward content-
adaptive switching of the receptive field at this scale, so the
mechanism's added freedom is unaligned with the task signal \\
\bottomrule
\end{tabularx}
\end{table}

\subsection{Mamba-2 novelty gate}\label{app:engaged-null:novelty}

The Mamba-2 novelty gate redeploys the LUCID kernel machinery from
value decorrelation to per-write novelty scoring, gating each token's
write into the SSD state via
$\omega_{t} = 1 / (1 + \gamma_{h} / (B_{t}^{\top} \Sigma_{c}^{-1} B_{t}))$
where $\Sigma_{c}$ is a chunk-local key-covariance estimate. The
mechanism is reported as a structurally inert null on the Mamba-2
spine, distinguished from the seven preceding entries by a
scale-mismatch failure mode rather than an axis-mismatch.

\begin{table}[h]
\centering
\small
\caption[Mamba-2 novelty gate scale-mismatch null]{Mamba-2 novelty
gate fact box.}
\label{tab:engaged-null:novelty}
\setlength{\tabcolsep}{6pt}
\renewcommand{\arraystretch}{1.05}
\begin{tabularx}{\linewidth}{@{}p{3.3cm} X@{}}
\toprule
Target axis & 2 (per-write associative-memory gating) \\
Cells & \texttt{mamba2\_novelty\_gate} (trainable $\gamma$) and
\texttt{mamba2\_novelty\_fixed\_g05} ($\gamma = 0.5$ fixed) \\
Dev / Test CER & $0.1186$ / $0.1186$ and $0.1187$ / $0.1187$ \\
$\Delta$ test vs vanilla Mamba-2 & $-0.0006$ and $-0.0005$ (tied
within $\sigma$) \\
Mobility evidence & Trainable $\gamma$ ends at $\approx 0.011$; fixed
$\gamma = 0.5$ produces test CER differing by $0.0001$ from the
trainable run \\
Structurally-inert reading & The Mahalanobis quantity
$q^{2} = B_{t}^{\top} \Sigma_{c}^{-1} B_{t}$ takes values in
$[10^{4}, 10^{5}]$ on real tokens, so the gate
$1 / (1 + \gamma_{h} / q^{2})$ requires $\gamma_{h} \sim 10^{3}$ to
produce meaningful attenuation. The trained and fixed $\gamma$
values explored are $4$ to $5$ orders of magnitude below the
required engagement threshold; this is a scale mismatch between the
gate's free parameter and the quantity it operates on, distinct
from the axis-mismatch nulls of
\Cref{app:engaged-null:delta,app:engaged-null:cayley,app:engaged-null:m2rnn,app:engaged-null:t2,app:engaged-null:s9b,app:engaged-null:a1prime,app:engaged-null:cb3} \\
\bottomrule
\end{tabularx}
\end{table}

\section{Family-D quadratic-lift saturated cluster}\label{app:engaged-null:family-d}

A second class of closed-cell evidence is the saturated polynomial
value-lift cluster on the discovery-phase RWKV-6 spine. Four
parametrisations of the same axis-5 function class were tested with
matched depth and parameter budgets, and saturate to a narrow
performance band around dev CER $0.124$. The cluster confirms the
axis-5 saturation reading from a different angle than the
engaged-null catalogue: the function-class extension is exercised by
each parametrisation, but the LibriSpeech task does not reward the
content-adaptive feature-map richness on which the cluster operates.

\begin{table}[h]
\centering
\small
\caption[Family-D quadratic-lift cluster on RWKV-6
discovery]{Family-D quadratic-lift cluster, RWKV-6 discovery-phase
30-epoch budget. Each row is a different parametrisation of the
same axis-5 polynomial value-lift function class, with
matched depth and parameter budgets.}
\label{tab:results:family-d}
\setlength{\tabcolsep}{6pt}
\renewcommand{\arraystretch}{1.10}
\begin{tabularx}{\linewidth}{@{}X p{4.5cm} r r@{}}
\toprule
\textbf{Cell} & \textbf{Parametrisation} & \textbf{Dev CER} & \textbf{Test CER} \\
\midrule
\texttt{rwkv6\_hadamard\_n2}     & Diagonal $k \odot k$ feature lift  & 0.1253 & 0.1251 \\
\texttt{rwkv6\_qtail}            & Full $k \otimes k$ Kronecker, top-2 layers & 0.1260 & 0.1240 \\
\texttt{rwkv6\_qtail\_lowrank\_all} & Low-rank $k \otimes k$, all layers     & 0.1238 & 0.1240 \\
\texttt{rwkv6\_pom\_vlift}       & Polynomial value lift~\cite{picard2026pom} & 0.1254 & 0.1253 \\
\bottomrule
\end{tabularx}
\end{table}

The four parametrisations differ in algebraic surface (diagonal
$k \odot k$, full Kronecker, low-rank Kronecker, polynomial mixer
projector) but share the same underlying function class of
polynomial value lifts of order at least two. The corresponding
HLA streaming higher-order extension~\cite{zhang2025hla} is not run
on this spine but is mapped to the same axis-5 cluster in
Section~\ref{sec:related:long-content}. The cluster's saturation to a
narrow band is the converse evidence for the axis-5 component of the
synthesis chapter's invariant.

\section{Discussion}\label{app:engaged-null:discussion}

The seven engaged-null entries and the Family-D quadratic-lift
cluster jointly support the converse direction of the
mechanism-prior alignment law of
Chapter~\ref{ch:experiments}. Each entry demonstrates a mechanism
that genuinely extends the function class along a specific axis,
that engages under SGD as evidenced by parameter mobility, and that
fails to convert into measured CER reduction on the LibriSpeech axis-1
probe. The seven mechanisms target axes 2, 3, and 5 of the
five-axis decomposition; the cluster targets axis 5; together they
cover three of the five axes in the decomposition.

The reading is not that the listed mechanisms do not work. The
delta rule and the Cayley orthogonal transition are productive
mechanisms on tasks that exercise associative-memory capacity under
interference and on tasks that exercise non-Abelian state-tracking,
respectively; M$^{2}$RNN, the non-normal RSE family, and the readout
gauge each have a natural test bed where the corresponding axis is
exercised. The reading is that these mechanisms are not productive
on the LibriSpeech axis-1 probe, because LibriSpeech's task structure
does not reward the function-class extension each one provides. This
is the cross-experiment invariant the synthesis chapter draws on,
and the closed-cell evidence here is the converse half of the
two-sided argument: the productive mechanisms in the main matrix
demonstrate the predictive law in one direction, and the engaged-
null catalogue demonstrates it in the other.

The Mamba-2 novelty gate is reported alongside the seven
engaged-null entries with a distinct reading. The novelty gate's
null is not an axis mismatch but a scale mismatch: the gate's free
parameter does not span the operating range of the quantity it
operates on, so the mechanism is structurally inert at the
architecture's regime regardless of the task's structural prior.
The distinction matters because it isolates the engaged-null
pattern from the broader category of unhelpful mechanisms: the
seven engaged-null entries proper require the function-class
extension to be exercised by parameter mobility before the null
becomes evidence for axis mismatch, and the novelty gate's
operational disconnection violates that precondition. Reporting the
two patterns separately preserves the precision of the
cross-experiment invariant claim.
